\section{선행연구 검토}

\subsection{Nowcasting 연구}

Nowcasting은 공식 통계 발표 전에 고빈도 지표를 활용하여 현재 시점의 거시경제 변수를 추정하는 기법임 \cite{banbura2012nowcasting}. Bańbura 등 (2012)은 혼합 빈도 동적 요인 모형을 활용한 nowcasting 프레임워크를 제시하였으며, Bok 등 (2017, 2019)은 FRBNY Staff Nowcast의 방법론을 상세히 설명함 \cite{bok2017macroeconomic, bok2019frbny}. COVID-19 팬데믹 기간 동안 Schorfheide와 Song (2020)은 혼합 빈도 VAR의 예측 성능 저하를 관찰하며, 비모수적 접근법의 효과성을 제시함 \cite{schorfheide2020nowcasting}.

\subsection{동적 요인 모형}

Stock과 Watson (2002)은 주성분 분석을 활용한 요인 추출이 고차원 시계열 데이터의 차원 축소에 효과적임을 보여줌 \cite{stock2002forecasting}. 최근 Andreini 등 (2020)은 자기인코더를 활용한 심층 동적 요인 모형(DDFM)을 제안하여 비선형 요인 구조를 학습할 수 있음을 보여줌 \cite{andreini2020deep}. Kim (2024)은 한국 데이터에서 딥러닝 기반 모형이 DFM보다 우수한 성능을 보인다고 보고함 \cite{kim2024deep}.

\subsection{연구 공백 및 기여}

기존 연구들은 다양한 모형을 체계적으로 비교한 연구가 부족하며, 특히 한국 거시경제 데이터에 대한 포괄적인 비교 연구가 제한적임. 본 연구는 이러한 공백을 메우기 위해 4개 모형을 동일한 데이터셋과 평가 기준으로 비교하는 프레임워크를 구축함.
